% THETA3-B24 contribution by /root/definite_bridge, 21 November 2026.
\subsection{THETA3-B24: the actual fixed-kernel inverse for a definite middle source}\label{B24}

In the totally definite natural-source sector, the integral
Gram matrix of B17 becomes a finite algebraic Gram matrix
times one common comparison scalar.  This identifies the
actual fixed-kernel inverse projectively with an algebraic
correspondence.  The same-type condition in A7 is retained:
this sector by itself gives divisor-generated balanced
tensors, rather than the general exceptional balanced type.

\paragraph{Fixed data and the scalar source type.}
Use the rank-three target $V$ and rank-two multiplicity
space $P=W_1\perp W_2$ of A7, with $d_1=d_2=d$ as
complete sign functions.  Then $P$ is totally definite.
Fix the same additive character, Weil splittings, target
and source levels, coefficient frames, and polynomial
Gaussian as in B17.  Require algebraic finite Schwartz
coefficients and algebraic values of the finite splitting
characters on the finitely many translates used below.
These conditions hold for the usual finite Weil model
with algebraic splitting characters.  Give
$[S]=[\U(P)]$ probability Haar measure and retain the
target conversion constant $\nu_G$ of B23.

Let $K(g,h)$ be the actual canonical scalar coefficient
kernel, including the determinant cocycle.  There is one
unitary character $\kappa_S$ of the full compact group
$S_\infty$ for which
\[
                 K(g,hk)=\kappa_S(k)K(g,h)
                    \quad(k\in S_\infty).
 \tag{B24a}
\]
This is stronger than just finite maximal-compact type.
Indeed the source action preserves the positive
Lagrangian used by the Gaussian at every real place.
At the active place its degree-two exterior factor is
$\bigwedge^2P_\mu^*$, a line; the half-determinant
factor is also a line.  At every other place, equality
of the two signs makes the chosen whole embedding
summand source-stable, and the Gaussian is a vacuum
line.  Their product transforms by a character.
These are precisely the determinant and exterior
factors in 032, Lemma 4.1 and equations (4.6)--(4.8),
printed pp.15--16.
The specified Weil splitting transports the bare
half-determinant characters to this honest character
of $S_\infty$.  This proof does not equate bare and
circle-valued splitting conventions.  Compactness
makes the character unitary.  It is independent of
the torus characters subsequently integrated in A7.

Choose finite representatives $s_b$ for
$S(F)\backslash S(\A_{F,f})/K_{S,f}$.  Let $D_S$ be
the diagonal probability-mass matrix.  Its entries
are positive rational numbers proportional to the
reciprocals of the finite stabilizer orders.
An incompatible stabilizer character forces the
corresponding scalar row to vanish; remove that row.
The surviving masses are not renormalized after
this removal.
The effective integrating source module has right
infinity type $\kappa_S^{-1}$, since B17 uses the
bilinear integral $Tu=\int K(g,h)u(h)\,dh$.
It is determined by its values at the $s_b$, and
its Hermitian norm is $z^*D_Sz$.
The operator on unrestricted $L^2([S])$ factors
through this finite module: averaging over
$S_\infty$ and $K_{S,f}$ annihilates the other
isotypic components.  Thus this reduction does not
alter B17's actual operator $TT^*$.

\paragraph{Algebraicity of the finitely many translated inputs.}
At a finite place, the Schr\"odinger Weil action is
generated by linear substitutions, additive-character
multiplications, and Fourier transforms.  Applied to
a locally constant compactly supported function,
each Fourier transform is a finite sum of roots of
unity times lattice volumes; the determinant factors
are powers of a rational residue cardinality, with
specified square roots, and the finite Weil indices
are roots of unity.  Rational or self-dual Haar
normalizations give algebraic lattice volumes.
Thus all these numbers are algebraic.
For explicit generator formulas with the splitting
character retained, see Chao Li, \emph{Geometric and
arithmetic theta correspondences}, Section 2.1,
printed pp.2--3: the Levi factor, unipotent factor,
and Fourier transform have exactly these forms.
The same argument applies after restricting the
ambient symplectic Weil model to the fixed unitary
splittings; their additional finite character
values are the hypothesis already stated.
The splitting-character values have been retained
as an explicit hypothesis above.  It follows that
the finitely many source and target finite translates
of the chosen input are again Schwartz functions
with algebraic values.  A common level and lattice
contain all of them.

Each such finite function is a finite algebraic sum
of tensor products for $W_1\perp W_2$: choose product
cosets fine enough to support and trivialize the
function.  The canonical determinant coefficient
therefore is a finite algebraic sum of wedges of
the two same-type theta one-forms.  Corollary 6.7
of 032 applies to this finite collection.  After
one common principal cover, all its kernels lie
in the cup image of a finite product of the actual
CM sources, as in B23.

They lie in one simple CM type with multiplicity,
and in one of its $(2,0)$ embedding components.
Here is a precise way to see this assertion even
if the original CM type is imprimitive.  In the
two copies of $U(v)$, project their tensor product
to equal $E$-embedding labels.  The rational
idempotent is supplied by the separable algebra
$E\otimes_\Q E$; its image is a rank-one
$E$-Hodge structure whose type at each embedding
is $(2,0)$ or $(0,2)$ according to $v$.  Its
Mumford--Tate characters form a single Galois
orbit, possibly with repetitions, because the
embeddings of the field $E$ form a single orbit.
It is consequently isotypic as a rational Hodge
structure.  The chosen embedding line belongs
to one embedding component of that simple type.
Every rational pair of CM source maps followed
by cup product preserves this assertion.
The complex linear spans in Corollary 6.7 do
not change the component.  Thus the finite
kernels belong to a full block
$\mathcal E_\chi$ of the kind in B23; their
image can be a proper subspace of that block.

\paragraph{One common CM evaluation scalar determines the source matrix.}
Choose the rational negative line and frame before
fixing the kernel, as in B17.  Select the finite
target translates of A5 separating the full block
$\mathcal E_\chi$, and use its coherent basis
$\omega=P_\chi\beta$.  Write $A$ for B23's
algebraic full-column-rank evaluation matrix, so
the actual target evaluation matrix is $t_xA$.

B19 applies at this one fixed CM period point
to all the finitely many translated inputs.
Choose one common finite theta module and its
algebraic Heisenberg isomorphism for each line.
Its determinant and tangent factors give one
nonzero constant $C_x$ for all these values.
With the oscillator constants, level factors,
and tangent matrices kept exactly as in B19,
there is a matrix $B\in M_{r_0\times b_0}(\overline\Q)$
such that
\[
       K(x_\ell,s_b)=C_x B_{\ell b}.
 \tag{B24b}
\]
Here $r_0$ is the number of target evaluation rows
and $b_0$ the number of surviving source cosets.
No torus integral is replaced by a finite sum:
(B24b) is a theta-kernel value.  The only
infinity integral used to reduce the source
has already been evaluated by (B24a).

Expand $K(\,\cdot\,,s_b)=\sum_a\omega_a E_{ab}$.
Choose a left inverse $L_A$ of $A$ over
$\overline\Q$, by row elimination.  Comparing
the two evaluations gives
\[
          E=\frac{C_x}{t_x}E_{\rm alg},
                  \qquad E_{\rm alg}=L_AB.
 \tag{B24c}
\]
The equality is exact because the evaluation
family separates the block.  This proves that
the image $\mathcal E_0=\operatorname{im}T$
is defined over $\overline\Q$ in the coherent
de Rham coordinates: it is the column span of
$E_{\rm alg}$.  It is not an assertion obtained
merely by declaring fibre evaluation to be a
Hodge operation.

Assume $\mathcal E_0\ne0$.  Choose a full-rank
algebraic column matrix $F$ spanning it and write
$E_{\rm alg}=F E_0$.  In the basis $\omega F$ put
\[
       G_0=F^*G_BF,
       \qquad W_0=E_0D_SE_0^*.
 \tag{B24d}
\]
Both are positive definite algebraic Hermitian
matrices.  Positivity of $W_0$ follows because
$E_0$ has full row rank.  On this basis the
bilinear integral $T$ is the matrix
$(C_x/t_x)E_0D_S$.  Its source Gram matrix is
$D_S$, and its target Gram matrix is
$\nu_G|P_\chi|^2G_0$.  Direct calculation of
the adjoint therefore proves
\[
 \begin{split}
 Q_T=TT^*&=
      \nu_G|P_\chi|^2\left|\frac{C_x}{t_x}\right|^2
                                      W_0G_0,\\
 Q_T^{-1}&=
      \frac{|t_x|^2}{\nu_G|P_\chi|^2|C_x|^2}
                             G_0^{-1}W_0^{-1}.
 \end{split}
 \tag{B24e}
\]
In particular this is the actual B17 inverse,
not the different finite restriction inverse
of A5.

The last matrix is realized by an algebraic
correspondence on the CM cup image.  Indeed,
the $G_B$-orthogonal projector onto the
column space of $F$ is the algebraic matrix
$F(F^*G_BF)^{-1}F^*G_B$.  Extending the
operator on that space by zero gives an
algebraic matrix on the full multiplicity
block.  B23's CM-source retract realizes
every such matrix by a correspondence.
The use of the transpose and the inverse source
polarization here is precisely B23's construction,
with the supplied Grothendieck \emph{cote 44},
scans 7--8, and \emph{cote 16}, scans 43--46,
as the framework locators.  It does not invoke
an inverse-polarization assertion for all of $X$.
It is positive self-adjoint on the selected
complex image, because it is a positive
multiple of $Q_T^{-1}$.

Consequently B17's exact adjoint double lift
and its literal rank-zero \texttt{LastLaw}
and \texttt{SeeLast} specialization use one
algebraic correspondence and the one positive
scalar displayed in (B24e), uniformly for
all torus characters integrated against this
fixed kernel.  The original scalar $C_x$
retains the exact B19 normalization; B22 can
normalize and bound its analytic reference
factors.  Algebraic denominators, conjugate
sizes, and the scalar $P_\chi/t_x$ are not
discarded.  The finite positive matrix $W_0$
admits the lattice/determinant certificates
of B7 whenever their stated arithmetic data
are provided.

\paragraph{The extension to other finite compact types.}
For a totally definite source and a finite
sum of compact infinity types, the same
argument applies whenever all coefficient
values share the common comparison scalar
used in (B24b).  One first expands in
matrix coefficients of the compact
representations, in compatible algebraic
bases.  An algebraic positive invariant Hermitian
Gram matrix can be chosen: infinitesimal invariance
for an algebraic basis of the compact real Lie
algebra, together with any component representatives,
is a finite real linear system over the real
algebraic numbers.  For the connected unitary groups
here the Lie algebra equations suffice.  Averaging
any positive real Gram matrix gives a positive real
solution.  Algebraic points of the solution space
are dense, and positivity is open, so a positive
algebraic solution exists.  Algebraic Cholesky
factorization then gives an orthonormal algebraic
change of basis.  Haar probability integration is
algebraic in those bases: averaging an
operator between two irreducibles gives
zero for distinct types and a scalar
identity for one type; taking the trace
gives the scalar $1/\dim\rho$ in the
usual matrix-coefficient orthogonality
formula.  Finite stabilizer projections
are averages over the stabilizers.
Their orders, the dimensions, and the
chosen algebraic representation matrices
therefore give an algebraic source Gram.
This proves the extension under the stated
common-scalar hypothesis.  That hypothesis
is proved above for the actual scalar
determinant cocycle, not asserted for
arbitrary unrelated infinity types.

\paragraph{The remaining indefinite-source matrix.}
For general anisotropic $P$ with a
noncompact real factor, write a finite-rank
kernel as $K(g,h)=\sum_a\omega_a(g)e_a(h)$.
Its source Gram is
\[
 W_{ab}=\int_{[S]}e_a(h)\overline{e_b(h)}\,dh,
 \qquad Q_T=\nu_G|P_\chi|^2W G_B
 \tag{B24f}
\]
when its image is the specified full block.
There is no compact Schur reduction to
the finitely many values in (B24b).
Projective algebraicity of this positive
matrix $W$ is the precise remaining
condition for a projectively algebraic
inverse in these coordinates.  It cannot
be forced by an invertible algebraic
change of target basis: if $W$ is changed
to $CWC^*$ with $C$ algebraic, one such
matrix is a positive scalar times an
algebraic matrix if and only if the
other is.

Finally, A7 shows that this natural
totally definite rank-two source has
equal complete types.  A balanced second
pair then has the same equal types, so
the selected four-factor line is a
product of divisor lines.  Formula
(B24e) does not settle the indefinite
source in a general exceptional switch.
Nor does a finite sum of identities
for different spaces $P$ constitute
one \texttt{MainX} tuple: the recovered
\texttt{Comm} fixes its Hermitian spaces
and requires its upper spaces to be
\texttt{DefSp}.  The exact clauses are
\texttt{Development/AbelianVarietyTheory/Definitions/}
\texttt{LatticeFrames/ValLastX.lean},
\texttt{Comm}, lines 30--69, and
\texttt{LawsX}, lines 212--221.
